\section{图结构把联合分布拆成局部因子}
\label{sec:13-Machine-Learning/07-Graphical-and-Sequential-Models/01}

假设你是一个急诊科医生，面前有十个二元症状变量：发热、咳嗽、胸痛、呼吸困难、恶心……每个病人就是这十个位置上的一组 0/1。如果什么结构都不假设，描述十个症状的联合分布需要 \(2^{10}-1=1023\) 个独立数字。十个症状而已，已经要记一千多个概率。三十个症状就是十亿量级。

可你真的需要把每种可能的症状组合都单独记录一遍吗？一个同时发烧和咳嗽的病人，与一个只有发烧的病人，两者的区别不是你凭空编出来的——它受人体病理通路约束。给定病人得了流感，发烧和咳嗽的概率就有了结构；再给定体温读数，许多症状之间的关联可能直接消失。条件独立是真实世界给你的压缩工具，图只是把这种压缩画成了可视的结构。

\subsection{遭遇指数墙：联合分布无法直接存储}

先从最笨的方法出发。十个二元变量，写出完整的联合概率表需要 \(1023\) 个数。用链式法则展开：

\[
p(x_1,\dots,x_{10})
=p(x_1)p(x_2\mid x_1)p(x_3\mid x_1,x_2)\cdots
p(x_{10}\mid x_1,\dots,x_9).
\]

项数是十项，似乎比一千多个数好多了。但最后一项 \(p(x_{10}\mid x_1,\dots,x_9)\) 仍然是一个条件于 \(2^9=512\) 种前置组合的概率表。链式法则只是把一个大表拆成了十个小表，这些小表加总的参数量一点没少——乘积展开不改变参数规模。

你还需要从数据里估计这些数字。如果某些前置组合在数据里只出现一两次甚至零次，对应的条件概率就无从估计。数据永远稀疏，组合永远爆炸。

\subsection{Bayesian network：声明``给定父母后，别的我都不需要''}

图模型做了一件链式法则没做的事：它声明某些条件依赖不存在。设有向无环图（DAG）中，每个节点 \(X_i\) 只直接依赖它的父节点集合 \(\operatorname{pa}(i)\)。联合分布被分解为

\[
p(x_1,\ldots,x_d)
=\prod_{i=1}^{d}
p\bigl(x_i\mid x_{\operatorname{pa}(i)}\bigr).
\]

这和普通链式法则的区别在于：链式法则让 \(X_i\) 条件于排序中所有更早的变量；图分解只保留父节点。那些被略去的变量，图的拓扑结构替你声明了条件独立——

\[
X_i\perp
\text{非后代且非父节点}
\mid X_{\operatorname{pa}(i)}.
\]

如果每个节点最多有三个父节点，每个条件概率表只需要 \(2^3=8\) 个数字；十个节点总共可能只需要几十个数。不是你用技巧压缩了分布，是你用领域知识限制了分布可以是什么——在一开始就把不可能的结构排除在外。

但缺边的含义来自整个分解和条件化集合。不能看见图上两个节点之间没有边，就说它们始终不相关。它们可能在给定第三个变量时变得相关，也可能在给定第三个变量后变得独立。边的存在和缺失，是相对于整个条件独立声明体系才有意义。

\subsection{三种三节点结构：同样的变量，完全不同的条件化行为}

理解图对条件独立的编码，只需要嚼透三种三节点格局。假设三个二元变量，分别称为 \(X,Z,Y\)。

第一种是链：

\[
X\longrightarrow Z\longrightarrow Y
\]

知识经过 \(Z\) 传递：\(X\) 影响 \(Z\)，\(Z\) 再影响 \(Y\)。\(X\) 和 \(Y\) 在无条件时可能相关——如果知道 \(X\) 发生了，\(Z\) 的分布会偏移，进而 \(Y\) 的分布也跟着动。但一旦你知道了 \(Z\) 的确切值，\(X\) 再变动也不会给 \(Y\) 带来任何额外信息。给定中间节点，路就断了。

第二种是共同原因：

\[
X\longleftarrow Z\longrightarrow Y
\]

\(Z\) 同时影响 \(X\) 和 \(Y\)。不观察 \(Z\) 时，\(X\) 和 \(Y\) 可能通过 \(Z\) 共享变化而相关。知道 \(Z\) 以后，\(X\) 和 \(Y\) 各自变成了依赖 \(Z\) 的独立噪声，相关性消失。给定共同父节点，路径也断了。这就是流行病学里``控制混杂变量''的图形原理：你要给定那个同时影响暴露和结局的因素，才能看清暴露和结局之间有没有剩下关联。

第三种是碰撞结构：

\[
X\longrightarrow Z\longleftarrow Y
\]

\(X\) 和 \(Y\) 分别独自影响 \(Z\)。若 \(X\) 和 \(Y\) 先验独立，不观察 \(Z\) 时路径被碰撞节点阻断——两个独立原因各走各的。一旦你给定 \(Z\) 的值或其后代的值，\(X\) 和 \(Y\) 反而产生了关联。

举一个具体的例子。假设 \(X\) 是``患者有心脏病''，\(Y\) 是``患者摔断了腿''，\(Z\) 是``患者进了急诊室''。心脏病和摔断腿在一般人群中基本独立。但如果你只在急诊室的人群里做统计——相当于给定了 \(Z\)——就会发现心脏病和摔断腿的发病率负相关：进急诊室的原因要么是前者、要么是后者，知道患者是摔断腿进来的，就降低了他是心脏病发作进来的概率。两个独立原因在被共同效应筛选后呈现出虚假关联。Selection bias 的图形根源就是这个碰撞结构。看到分层后两个变量相关，不能自动断言其中一个造成了另一个。

d-separation 就是把这三条局部规则推广到任意大小的 DAG：给定条件集合 \(C\) 后，检查变量间的每条无向路径。若路径上每个节点要么是非碰撞节点且被 \(C\) 包含（路径阻断），要么是未被观察且未观察到后代的碰撞节点（路径也阻断），那么对应条件独立就有图分解的保证。

\subsection{Markov random field：当方向不重要时用势函数}

有向图天然适合表达生成过程：先有原因，后有结果。但有些场景里变量之间没有明显的先后方向——比如图像里相邻像素的标签，你只是希望相邻像素大概率取一致的类别，不存在谁生成谁。这时用无向图更自然。

无向图用势函数而非条件概率建模。若图中最大团集合为 \(\mathcal C\)，分布写成

\[
p(x)
=\frac1{Z}
\prod_{C\in\mathcal C}\psi_C(x_C),
\qquad \psi_C\ge0,
\]

\[
Z=\sum_x\prod_{C\in\mathcal C}\psi_C(x_C).
\]

势函数 \(\psi_C\) 不是概率——它不要求归一，只负责给团内各种配置打出相对分数。\(Z\) 是配分函数，把所有人的分数加总后，才能把势函数翻译成概率。

这里藏着一个代价：势函数是局部的，每个只关心团内几个变量，但配分函数 \(Z\) 是对所有可能配置全局求和。一个小团局部调参，要牵动全空间重算 \(Z\)，因为所有配置的概率都必须保持加起来是一。参数学习时

\[
\nabla_\theta\log Z(\theta)
=\E_{p_\theta}[T(X)],
\]

右手边是模型分布下的期望——你选的这个看起来只是局部的模型，在学习时又把全局期望拉了回来。MCMC、变分推断和对比近似这些技术，正是为了帮你绕开这个全局求和的精确计算。

无向图的局部 Markov 性质很干净：给定一个节点的所有邻居后，它和图中其他所有节点条件独立。与有向图不同，无向路径只需检查是否经过条件集合——没有碰撞节点的反向开启问题，规则更简单。

\subsection{同一套条件独立可以画出不同的图}

你现在知道图编码了条件独立。但反过来：给定一组条件独立声明，能唯一画出一张图吗？

不能。不同的 DAG 可以蕴含完全相同的条件独立集合——它们属于同一个 Markov 等价类。两个变量如果总是同增同减，因果方向、共同原因和共同效应这三种图，在纯观察数据面前是分不开的。

因此 Bayesian network 里的箭头首先表示的是概率分解的方向，而不是因果干预的方向。你只能说``\(X\) 在条件概率表的自变量列表里''，不能说``改变 \(X\) 会以图中箭头的方式引起 \(Y\) 的变化''。把箭头升级为因果解释，需要额外的结构方程、faithfulness 假设、无隐藏混杂、介入实验或工具变量等条件。一个预测表现出色的图模型，箭头方向可能和真实因果恰好相反——预测好不等于因果真。

\subsection{推断就是消灭乘积里的多余变量}

建模写完分解式，接下来最常做的事是：给定部分变量的观测值 \(x_E\)，计算另一些查询变量 \(x_Q\) 的后验分布。公式写下来就是多重求和：

\[
p(x_Q\mid x_E)
\propto
\sum_{x_H}
\prod_a\psi_a(x_a).
\]

变量消元法依次挑出包含某个待消变量的所有因子，相乘，再把该变量加掉。每条边都可能影响中间因子在半路上的大小，不同消元顺序产生的最大中间因子尺寸可以差好几个数量级。这个最大尺寸由图的 treewidth 控制，不单纯由节点数决定。

链和树上消息可以在与节点数成线性关系的时间内完成——HMM 的 forward--backward、树上的 sum-product、Kalman filter 的预测-更新，都是在用图结构的稀疏性反复执行局部消元。稠密图或者有大环的图，中间因子会指数增长成无法存储的大表。

这正是图模型的深层统一性：一条边的存在，既声明了条件依赖，又可能堵住一个消元顺序让它变贵。模型的统计假设和计算代价画在同一张图上——你看得见哪些独立帮你省钱，也看得见哪些依赖让你付账。

\subsection{学好一张图不等于拍到了现实的照片}

有限样本下，接近独立和真正独立很难区分。噪声让似然面变得平坦，许多候选图得分接近，选出来的所谓``最优图''可能只是这次抽样赢了一点点。隐藏变量的存在更麻烦——图上本来没有边的两个观测变量，若共享一个未观测的共同原因，数据里会呈现出残留关联，学出的图会凭空多出边来。

正则化、评分准则（BIC、BDeu 等）和搜索空间的限制都参与了结构选择。得到一张稀疏图不等于缺失边在现实中绝对没有作用——可能只是作用太弱或数据太少。得到的连边也同样不自动对应直接因果通道。

用图模型时应分别说清楚四层条件：图声明了哪些条件独立，参数规定了哪些局部分布的形态，推断算法精确或近似地在算什么，以及你手中的数据量是否足以把这些结构与随机波动区分开。图长得清爽不等于四层条件都站得住。Machine Learning Part 末尾的\hyperref[ch:120]{``因果推断：预测世界不等于改变世界''}会进一步为箭头填入结构方程和干预语义，并回答那些让统计学家夜不能寐的问题——哪些因果量能从观察分布中识别出来，哪些永远不能。
